\section{freestanding C++20 与跨语言 ABI}

C++ 是一种语言，不等于桌面操作系统提供的运行环境。内核可以使用类型系统、
封装、\code{constexpr} 和强枚举，但必须自行建立 ABI、内存、对象生命周期
与失败策略。

\topic{一、托管环境与独立环境}
C++ 标准区分 hosted 与 freestanding 实现。普通应用启动前，操作系统装载器
映射程序和动态库，语言运行时准备栈、线程局部存储、全局构造、堆和退出流程，
随后才调用 \code{main}。程序还可以把文件、时钟、终端和线程交给系统调用。

固件与内核本身正是这些服务的提供者。处理器跳入内核入口时，不存在另一个
操作系统替它调用构造函数或处理异常。目标使用
\code{--target=x86_64-unknown-none-elf} 和 \code{-ffreestanding}，
明确告诉编译器当前程序不依赖宿主平台。

\begin{osgridlongtable}{L{0.23\textwidth}L{0.29\textwidth}L{0.32\textwidth}}
\textbf{能力} & \textbf{普通应用中的来源} & \textbf{本项目中的建立方式}\\ \hline
\endfirsthead
\textbf{能力} & \textbf{普通应用中的来源} & \textbf{本项目中的建立方式}\\ \hline
\endhead
初始栈 & 进程装载器与内核 & 启动汇编选择保留 RAM 并对齐\\ \hline
零初始化区 & 装载器或 C 运行时 & 入口按链接符号主动清零\\ \hline
全局构造 & CRT 初始化数组 & 后续自研初始化遍历或避免依赖\\ \hline
动态内存 & libc 与系统调用 & 物理页、虚拟内存和内核堆\\ \hline
异常展开 & ABI 运行库与展开表 & 初期禁用，使用显式状态\\ \hline
程序退出 & 运行库调用操作系统 & 内核停机、重启或结束任务\\ \hline
\end{osgridlongtable}

\topic{二、编译器选项固定生成代码边界}
\code{-fno-exceptions} 避免生成异常展开和 \code{__cxa\_*} 依赖；
\code{-fno-rtti} 避免类型信息与动态转换运行时；
\code{-fno-stack-protector} 避免尚未提供的栈保护符号；
\code{-fno-threadsafe-statics} 和 \code{-fno-use-cxa-atexit} 避免局部静态
初始化锁与退出析构注册。\code{-fno-builtin} 阻止编译器在没有实现契约时
擅自把循环替换为 \code{memcpy} 或 \code{memset} 调用。

\code{-mno-red-zone} 对内核尤其重要。System V AMD64 ABI 允许用户态函数
使用 RSP 以下 128 字节而不移动栈指针，但异步中断可能在当前栈上压入现场，
破坏这片区域。禁用 red zone 后，编译器不会把临时值放在那里。

初期还禁用 MMX、SSE 与 SSE2。它们并非永远不能使用，而是操作系统尚未保存
这些扩展寄存器，也未建立每线程浮点状态。只有上下文切换协议覆盖相应状态后，
编译器才能安全生成这些指令。

\topic{三、ABI 是汇编与 C++ 的共同契约}
应用二进制接口规定参数寄存器、返回值、调用者与被调用者保存寄存器、栈方向
和对齐。它使分别编译的函数能够协作，也使启动汇编能够调用 C++ 入口。函数
签名只描述源语言类型，真正的跨边界行为还包括寄存器和内存约束。

64 位内核入口将采用 System V AMD64 调用约定的受控子集。跳转到 C++ 前，
汇编层要保证 RSP 对齐、DF 清零、必要寄存器携带 BootInfo 地址，并确保页表
覆盖代码、栈和参数。C++ 返回后的行为也要预先定义，通常进入 panic 或停机，
而不是落入未知字节。

\topic{四、精确宽度与领域类型}
裸机代码面对寄存器、页表项和磁盘字段，不能让 \code{int}、\code{long} 的
平台差异进入协议。项目采用 \code{uint8_t}、\code{uint16_t}、
\code{uint32_t} 和 \code{uint64_t} 等显式宽度类型。外部格式严格服从协议；
内部地址、容量和计数在没有更窄限制时优先使用 64 位，以避免后续扩展重新修改
整个接口。

固定宽度仍不足以区分物理地址、虚拟地址和字节数量。基础模块用
\code{PhysicalAddress} 与 \code{ByteCount} 包装同为 64 位的值，让函数
签名携带领域含义。显式构造函数阻止普通整数无意转换，\code{[[nodiscard]]}
提醒调用者不能丢弃失败状态，\code{noexcept} 固定当前没有异常展开的契约。

\topic{五、半开区间消除边界歧义}
地址范围表示为 \([begin,end)\)：起点属于范围，终点不属于范围。大小直接为
\(end-begin\)，相邻范围 \([a,b)\) 与 \([b,c)\) 不重叠，空范围满足
\(begin=end\)。这种约定与数组迭代器、页范围和 ELF 文件区间保持一致。

创建范围时不能直接计算 \code{begin + size} 再检查结果，因为无符号加法溢出
后已经回绕。实现先计算从起点到最大值的剩余空间：
\[
  \mathrm{maximum_size}=UINT64\_MAX-begin.
\]
实现只在下列判定条件成立时执行加法：
\begin{center}
  \code{size <= maximum_size}.
\end{center}
失败通过强枚举
\code{AddressRangeCreationStatus} 返回，不依赖异常、全局错误码或宿主
运行库。

\topic{六、头文件与源文件承担不同责任}
公共头文件描述其他模块可以依赖的类型、常量和函数契约；源文件保存实现细节。
调用者不需要看到成员函数怎样比较原始数值，修改实现也不应迫使所有上层代码
重新理解内部算法。模板和必须内联的定义是语言规则要求的例外，仍要放入明确的
\code{.tpp} 或头文件结构。

下面两份清单来自配套工程。头文件先建立领域语言，源文件再实现溢出检查和区间
关系。类成员函数访问成员时显式使用 \code{this->}，使成员来源在系统代码的
复杂表达式中保持清晰。

\lstinputlisting[
  language=C++,
  caption={地址与范围的公共类型}
]{../../../../source/foundation/include/os/foundation/address_range.hpp}

\lstinputlisting[
  language=C++,
  caption={地址范围的独立实现}
]{../../../../source/foundation/src/address_range.cpp}

\topic{七、常量替代魔法数字与魔法字符串}
硬件端口、位掩码、页大小和协议标记都要用
\code{OS\_项目模块\_功能} 形式的全大写常量表达。名称使评审者能够追溯数值
来源，也让同一硬件事实只维护一次。优先使用有类型的 \code{constexpr} 或
\code{const} 对象，只有汇编器、链接器或条件编译确实需要文本替换时才使用宏。

“没有字面量”并不意味着把 \code{0} 和 \code{1} 机械隐藏到模糊名字中。
项目目前对它们也采用领域常量，是为了统一强调地址计算的语义。更重要的规则是
任何协议值、容量上限和特殊标记都必须有准确来源，不能把知识埋在表达式里。

\topic{八、volatile、原子操作与内存顺序}
\code{volatile} 表示一次访问具有必须保留的外部可见性，适合按规范访问某些
MMIO 寄存器，但它不提供线程同步、原子性或处理器内存屏障。端口映射 I/O 使用
\code{in}/\code{out} 汇编边界，同样要明确编译器约束和设备顺序。

多核同步要使用原子指令、内存顺序和锁协议；设备 DMA 还涉及缓存一致性与 I/O
屏障。把共享变量声明为 volatile 不能修复竞态。当前单核、关闭中断的复位阶段
尚未进入并发模型，后续引入中断和多核时再建立这些机制。

\topic{九、现代语言特性服务于可证明约束}
\code{enum class} 防止不同状态域意外混用；\code{constexpr} 允许在编译期
验证布局和掩码；概念与模板可用于受约束的寄存器或容器抽象；
\code{std::span} 一类无所有权视图适合表达指针加长度，但前提是项目提供所需
freestanding 标准库能力。现代并不等于默认启用一切标准库组件。

采用一个特性前要检查生成代码、运行时符号、异常行为与二进制大小。代码保持
接近领域语言，同时每个抽象都能下降到清楚的地址、寄存器和控制流，这才是
“代码即文档”在内核中的含义。
